\documentclass[10pt]{article}
\usepackage[a4paper,margin=24mm]{geometry}
\usepackage{amsmath,amssymb,amsthm,booktabs,longtable}
\newtheorem{theorem}{Theorem}\newtheorem{lemma}{Lemma}
\newcommand{\Aut}{\operatorname{Aut}}\newcommand{\Out}{\operatorname{Out}}
\title{The printed almost-simple chiral atlas: complete applications and its domain}
\author{Private L070 recovery component}\date{}
\begin{document}\maketitle
The fixed source is Hartley--Hubard--Leemans, \emph{Two atlases of
abstract chiral polytopes for small groups}, Ars Math. Contemp. 5 (2012),
371--382. We distinguish its 79 printed group models, their 88 occurrences
in Tables 2--7, and the larger domain stated in Section 4. All computations
below use those exact named models; printed polytope counts are never
search filters. The separate first-atlas classification is not repeated here.

\paragraph{Convention and retained recognition.}
An admissible prefix marking is an ordered tuple
\[\alpha=(1,\alpha_1,\ldots,\alpha_{r-1})\]
generating $G$, with the
adjacent orders greater than two, nonadjacent differences involutions,
and the full intersection property for the subgroups generated by
$\alpha_i^{-1}\alpha_j$ on each index set. The acquired recognition
theorem gives the corresponding regular or chiral abstract polytope.
A marking is regular exactly when the simultaneous inversion of
$\alpha_1,\ldots,\alpha_{r-1}$ extends to an automorphism.
Markings are identified by ordered generator images under $\Aut(G)$;
reversal is not additionally identified. In the chiral part the source
counts the two enantiomorphic forms separately: Sections 3 and 4 give the
ordered-marking construction, and Section 6 expressly divides its count
by two only when those forms are to be identified. This convention was
independently checked in the acquired source-orientation review.

\begin{theorem}
Relative to the acquired named models, recognition and exhaustive-prefix
theorems, and the stated finite group-operation implementations, all
chiral columns for the 79 printed groups are determined. Their total is
\[
 16974=15056\text{ (rank 3)}+1890\text{ (rank 4)}+28\text{ (rank 5)}.
\]
Of the 88 printed rank distributions, 87 agree. Table 6, page 381,
prints 26 for $\mathrm{PSU}_4(2)$; the correct rank-three count in the
same convention is 46. There are no other rank-count discrepancies in
the printed chiral rows.
\end{theorem}
\begin{proof}
The accepted prefix theorem enumerates inner-conjugacy classes by
choosing one representative of each nonidentity, noninvolutory first
conjugacy class, and then representatives under the centralizer of the
prefix subgroup at every extension. Its full intersection test and
reflection test are unchanged. Every extension strictly increases that
subgroup, so the finite enumeration exhausts all ranks, rather than
stopping at the largest rank printed in the source.

The actual executions consist of 72 clean complete whole-group calls,
one complete 25-class union for $S_9$, and six complete fixed-plan unions
for $L_3(4).D_{12}$, $P\Sigma U_3(5)$, $J_2$, $L_3(5):2$,
$P\Gamma U_3(5)$ and $M_{22}:2$. The last six replace six earlier
nonreturns, which receive no census credit. Each fixed-plan call checks
its full pinned first-representative list against freshly computed
complete classes by actual conjugating transporters. The receipt union
requires the exact same literal list, model, engine and acceptance, and
each index exactly once. In particular an index from a different
randomly ordered class list cannot be substituted.

The application reader reconstructs the whole-call drivers and exact
input/output inventories, checks all ordered class-return markers and
the terminal census marker, and requires a clean terminal process. For
class unions it invokes the unchanged accepted v2 or v3 receipt verifier
in isolated Python processes. The fresh application makes no group calls.
All literal census files equal their previously acquired unions. The
older whole-call receipts bind their exact program, model and driver,
but record only the GAP executable locator, not a complete runtime-file
trace; that historical provenance limit is retained explicitly. The
stronger v2/v3 runtime inventories are not retroactively attributed to them.

Finally $\Out(G)$ acts freely on inner classes of generating markings.
Indeed, an automorphism stabilizing an inner class agrees on every
generator with an inner automorphism, hence equals it. Regularity, rank
and ordered type are invariant. Thus each returned rank/type/reflection
stratum has size divisible by $|\Out(G)|$, and division gives exactly the
ordered $\Aut(G)$ classes. The reader checks these divisions before
joining all 79 groups to all 88 source occurrences. The resulting complete
machine-readable row list gives the stated total and sole discrepancy.
The previously accepted independent literal 46-class certificate for
$\mathrm{PSU}_4(2)$ corroborates this row without assuming exhaustion or
outer-index division; it is acquired separately, not counted twice.
\end{proof}

\paragraph{The Section 6 aggregates.}
The accepted regular-column application gives 9210 regular classes, with
rank distribution $8319,810,61,16,3,1$ for ranks 3 through 8. Its only row
correction is $\mathrm{PSL}_2(27)$, 27 to 26. Combining that result with
the theorem gives the following family totals; overlapping group
occurrences remain in their individual families.
\begin{center}
\begin{tabular}{crrrr}\toprule
Source table&Regular&Chiral&Regular rank $\geq4$&Chiral rank $\geq4$\\\midrule
2&1536&4780&271&360\\
3&804&1710&234&348\\
4&3903&1882&45&338\\
5&1846&3998&180&448\\
6&1202&4242&170&464\\
7&14&412&0&0\\\bottomrule
\end{tabular}
\end{center}
The nine duplicate occurrences contribute 95 regular and 50 chiral
classes and are removed only in the global sum. The source's table-total
arithmetic errors, already visible without a census, are Table 3 chiral
$1702\to1710$, Table 4 chiral $2050\to1882$, and Table 5 regular
$1840\to1846$. The two newly verified row corrections then give the
global $9205\to9210$ regular and $17114\to16974$ chiral totals.
The Section 6 comparison for the $\mathrm{PSL}_2$ family has rank-at-least-four
counts 45 regular and 338 chiral, rather than 51 and 506.

Mirror reversal is a fixed-point-free involution on the chiral ordered
classes: a fixed class would provide the reflection automorphism excluded
by chirality. Therefore identifying enantiomorphic forms gives
$16974/2=8487$ classes, by rank $7528,945,14$. This does not identify
duals. Under this convention regular classes still outnumber mirror
pairs, whereas the two oriented forms counted separately outnumber
regular classes. The qualitative sporadic, unitary and Suzuki comparisons
in Section 6 remain true. Its older 5265 count is explicitly up to both
isomorphism and duality, and the Conder genus counts are attributed
historical context from other work; neither is inferred from this census.
The source excludes $\Aut(J_2)$ from its regular aggregate, as do our
79 printed models. $J_2:2$ has order 1209600 and lies outside both
numerical cutoffs discussed below.

\begin{lemma}[Automorphisms of an almost simple preimage]
Let $S$ be nonabelian finite simple and $S\leq G\leq\Aut(S)$.
Then restriction identifies
\[
 \Aut(G)\cong N_{\Aut(S)}(G),\qquad
 N_{\Aut(S)}(G)/S=N_{\Out(S)}(G/S).
\]
\end{lemma}
\begin{proof}
The centralizer of $S$ in $\Aut(S)$ is trivial. Any nontrivial normal
subgroup of $G$ contains $S$: otherwise its commutator with $S$ is
trivial, which places it in that centralizer. Hence $S$ is the unique
minimal normal subgroup and is characteristic. An automorphism of $G$
restricts to $\beta\in\Aut(S)$. Its value at $g\in G$, considered as an
automorphism of $S$, must be $\beta g\beta^{-1}$, by conjugation on $S$.
It therefore embeds in the displayed normalizer. Conversely a normalizer
element induces an automorphism of $G$ by conjugation. The same trivial
centralizer makes this action faithful. The last equality follows since
$G$ is the full preimage of $G/S$.
\end{proof}

\paragraph{The group-order and automorphism-name columns.}
The accompanying 79-model registry rejoins every actual ordered
generator array to its accepted named-model selector; natural alternating
and symmetric models are identified by permutation parity and their
checked full factorial order. These are identifications, not naming by
order alone. Two printed group orders are wrong: $M_{22}$ is 443520,
not 443510, and $\mathrm{PSL}_2(25).2$ is 15600, not 7800.

For automorphism names we retain the full named simple-group
automorphism classifications and the accepted outer-preimage
identifications. The classified-input locators and the exact named target
model for every row are recorded separately. Except for $L_3(4)$ and
$U_3(5)$ the relevant outer groups are abelian, so every outer subgroup
has full normalizer. In $D_{12}\cong\Out(L_3(4))$, the central
order-two subgroup is normal, each reflection subgroup has a four-group
normalizer, and that four-group is self-normalizing. The other relevant
subgroups are normal. The already accepted ATLAS-word joins distinguish
the central $2_1$, field $2_2$ and graph $2_3$ types. Thus precisely the
two reflection preimages and their four-group preimage have automorphism
group $L_3(4).2^2$; the others have the full $L_3(4).D_{12}$.
In $S_3\cong\Out(U_3(5))$ the reflection preimage
$P\Sigma U_3(5)$ is self-normalizing; the other three preimages have
full automorphism group $P\Gamma U_3(5)$. All these small normalizers
are computed literally in the accompanying reader.

The source's $\Aut(\mathrm{PSU}_3(3))=P\Gamma L_3(3)$ is a label
error: the retained unitary automorphism theorem gives $P\Gamma U_3(3)$,
of order 12096. In contrast, $P\Gamma L_3(3)=PGL_3(3)$ has order
\[
 \frac{(3^3-1)(3^3-3)(3^3-3^2)}{3-1}=\frac{11232}{2}=5616,
\]
since $\mathbb F_3$ has no nontrivial field automorphism. For $J_2$,
the acquired full action is the accepted pair
$\operatorname{PrimitiveGroup}(100,1)$ inside
$\operatorname{PrimitiveGroup}(100,2)$. A saved six-letter word in the
latter generators gives an actual outside involution; a clean bounded
verification checks the exact arrays, normalization, index two and
generation. Thus the printed colon $J_2:2$ has an actual complement,
not merely an outer quotient of order two. The unsuccessful initial
length-two proposal and the longer proposal's GAP parse warning are
preserved; only the clean fixed-witness verification supplies this join.

\begin{theorem}[Historical-domain boundary, with classified inputs]
Using the retained CTblLib transcription of the 1985 ATLAS contents and
the accepted complete simple/outer roster below one million, the
Section 4 socle condition gives 80 almost simple isomorphism classes
of order below one million. Exactly 79 have order below 900000, and
their named models are precisely the 79 printed groups. The additional
class is $\mathrm{Sp}_4(4)$ of order 979200.
\end{theorem}
\begin{proof}
The accepted L054 domain theorem gives every almost simple group below
one million as a preimage of an outer subgroup. Its unique-minimal-normal
argument proves that isomorphisms between such preimages are exactly
outer conjugacies. Its finite outer-subgroup lists and complete simple
roster are acquired, with their explicit classification premises.

The pinned CTblLib \texttt{AtlasPages} record has 97 entries. Its own
documentation identifies the original book's page-v contents and later
additions. Exactly 93 entries have a nonempty ordinary-ATLAS page list;
the other four are later additions $L_2(49),L_2(81),L_6(2),S_{10}(2)$.
Retaining only the original ordinary-page entries intersects the accepted
small-order roster in 32 socles and 80 almost simple classes. The literal
table identifiers for the 79 acquired models are distinct and are
exactly this list minus \texttt{sp44}. Their checked orders are all
below 900000. The remaining group has order 979200 and ordinary ATLAS
page 44. This proves both strict-cutoff assertions without guessing an
old website edition or using polytope counts to identify groups.
\end{proof}

\paragraph{The precise source limit.}
The abstract on page 371 says that $\Gamma$ has order less than 900000
and is listed in the ATLAS. Section 4, page 375, instead explicitly
uses a socle $S$ appearing there and $|\Gamma|<10^6$. Our theorem
implements that explicit socle condition. The 79 printed models agree
with its 900000 restriction, but not with its one-million restriction.
We have not inspected a scan of the 1985 contents: the pinned CTblLib
transcription is an openly retained premise. We do not assert that every
extension has its own separately printed 1985 table. The order boundary
and this $\Gamma$-versus-$S$ wording distinction are kept separate.
No polytope census for $\mathrm{Sp}_4(4)$ is supplied by this component;
no conclusion about an omitted zero count is licensed. Consequently this
application closes the printed second-atlas entries and their aggregate
corrections, but does not silently convert the larger Section 4 assertion
into the narrower printed roster.

\paragraph{Attribution and limits.}
Root authored the prefix algorithm/theory, most original whole-group
applications and the fixed queue, and suggested the normalizer lemma.
The present author wrote the v2/v3 receipt repair, ran the four assigned
fixed-plan queues, and assembled the present exact joins and small
normalizer calculations. The corresponding independent reviews are
acquired with these roles unchanged. The regular application and U4(2)
positive correction have their separate independent reviews. Publication
provided a bounded Section 6 source-only reading. No new polytope census,
whole-manuscript acceptance, open-question resolution, Lean proof or
compression measurement is asserted here.
\end{document}
